\section{Introduction}\label{sec:introduction}

The Grothendieck--Teichm\"uller Lie algebra records the infinitesimal
symmetries of the associativity and braiding constraints for
parenthesized braids.  Its defining equations are short, but they place
conditions in every weight on a free Lie algebra in two letters.
The Deligne--Drinfeld conjecture predicts that these conditions leave
exactly a free Lie algebra with one generator in each odd weight
starting at three.  We prove this conjecture, including generation of
the entire solution space.

\subsection{The statement}
Let \(L=\Lie_{\Q}\langle x,y\rangle\), graded by the total number of
letters.  We call this degree the \emph{weight}.  For a finite set of
labels, the infinitesimal pure braid Lie algebra \(\mathfrak t_n\)
has generators \(t_{ij}=t_{ji}\), \(i\ne j\), and relations
\begin{equation}\label{eq:braid-relations}
 [t_{ij},t_{k\ell}]=0\quad(i,j,k,\ell\text{ distinct}),\qquad
 [t_{ij},t_{ik}+t_{jk}]=0\quad(i,j,k\text{ distinct}).
\end{equation}
Every generator has weight one.
Let \(W\subset L\) be the graded vector space of polynomials satisfying
\begin{align}
 \psi(x,y)+\psi(y,x)&=0,\label{eq:antisymmetry}\\
 \psi(x,y)+\psi(y,-x-y)+\psi(-x-y,x)&=0.\label{eq:three-term}
\end{align}
\begin{equation}\label{eq:pentagon}
\begin{aligned}
 &\psi(t_{12},t_{23}+t_{24})+\psi(t_{13}+t_{23},t_{34})\\
 &\qquad=
 \psi(t_{23},t_{34})+\psi(t_{12}+t_{13},t_{24}+t_{34})
 +\psi(t_{12},t_{23}).
\end{aligned}
\end{equation}
The last equation is imposed in \(\mathfrak t_4\).
These are the usual defining equations of \(\grt\);
see \cite[Section~6.1]{Willwacher}.  We work first with their polynomial
solution space \(W=\bigoplus_n W_n\), where \(W_n\) is its weight-\(n\)
piece.  The usual completed algebra is \(\widehat W=\prod_n W_n\).

For \(\psi\in L\), define the derivation
\begin{equation}\label{eq:ihara-derivation}
 D_\psi(x)=0,\qquad D_\psi(y)=[y,\psi].
\end{equation}
The Ihara bracket, with the convention used throughout this paper, is
\begin{equation}\label{eq:ihara}
 \{\psi,\phi\}=D_\psi(\phi)-D_\phi(\psi)+[\psi,\phi].
\end{equation}

\begin{theorem}\label{thm:main}
There are homogeneous elements
\(\sigma_{2k+1}\in W_{2k+1}\), one for each \(k\ge1\), for which
\[
 \Lie_{\Q}\langle e_3,e_5,e_7,\ldots\rangle
 \longrightarrow (W,\{\, ,\,\}),\qquad
 e_{2k+1}\longmapsto\sigma_{2k+1},
\]
is an isomorphism of graded Lie algebras.
In particular, \(W\) is closed under the bracket
\eqref{eq:ihara}.  The induced map on completions by weight is a
continuous graded isomorphism.
\end{theorem}

The generators in Theorem~\ref{thm:main} are not canonical.
The assertion combines two requirements: there are no additional Lie
relations among them, and they span every solution of
\eqref{eq:antisymmetry}--\eqref{eq:pentagon}.

\subsection{Context and the main estimate}
The arithmetic study of the fundamental group of the projective line
minus three points provides a central setting for these symmetry
questions; Deligne developed its motivic and tangential-basepoint
framework in \cite[Introduction and Section~15]{Deligne}.
Drinfeld's associators express compatibility between reassociation and
braiding, and their symmetries led to the
Grothendieck--Teichm\"uller group and its Lie algebra
\cite{Drinfeld}.  The equations above are the corresponding
infinitesimal compatibility conditions.

The conjecture has three parts: existence of the proposed odd-weight
elements, freeness of the Lie algebra they generate, and exhaustion of
the equation space.  Drinfeld's Proposition~6.3 credits Ihara with the
odd-weight existence result and gives an associator proof.
Brown's work on mixed Tate motives over \(\mathbb Z\) proves the
faithfulness of the motivic action on the fundamental group of the
three-punctured line \cite{Brown}.  The graded Lie algebra of the
prounipotent motivic Galois group is free on generators corresponding to weights
\(3,5,7,\ldots\); its inclusion in the associator symmetry algebra
therefore supplies a free Lie subalgebra of \(W\)
\cite[Section~1.4.2]{BrownDepth}.
Willwacher subsequently proved freeness for any homogeneous family of
solutions, one in each odd weight \(2k+1\ge3\), whose corresponding
element has nonzero coefficient of
\(x^{2k}y\) \cite[Theorem~1.2]{WillwacherFree}.
This coefficient lies in the \emph{depth-one} component, consisting of
words with exactly one occurrence of \(y\).

Freeness leaves open whether there are additional solutions of the
defining equations; see \cite[Section~9]{Willwacher}.
Naef and Willwacher's computations of the linearized
Kashiwara--Vergne algebra, together with the known inclusions,
establish the required equality through weight~29
\cite[Theorem~5 and Corollary~7]{NaefWillwacher}.
Our main estimate is the all-weight upper bound
\[
 \dim_{\Q}W_n\leq
 \dim_{\Q}\Lie_{\Q}\langle e_3,e_5,e_7,\ldots\rangle_n,
 \qquad \wt(e_j)=j,
\]
proved in Corollary~\ref{cor:dimension-bound} for the full rational
equation space.  The known free subalgebra gives the reverse inequality.
The proof here also obtains that lower bound internally:
Sections~\ref{sec:category}--\ref{sec:induction} construct a rational
subspace of solutions closed under the Ihara bracket, supply its
odd-weight depth-one values, and choose integral generators attaining
the bound.  These choices connect the rational construction to the
characteristic-two argument.

The completed assertion has a graph-theoretic consequence.  Over
\(\Q\), Willwacher's isomorphism
\cite[Theorem~1.1]{Willwacher} identifies the zeroth cohomology of
Kontsevich's graph complex with \(\widehat W\).  Thus
\[
 H^0(\mathrm{GC}_2)\cong
 \widehat{\Lie_{\Q}\langle e_3,e_5,e_7,\ldots\rangle}.
\]
The graph complex and its completion are taken in the conventions of
that source.  The superscript zero is a cohomological degree, whereas
the generator indices specify weights; no other cohomological degree
is computed here.

\subsection{Proof strategy}

The main step is a degeneration of the pentagon in characteristic two.
Inside the free associative algebra, put
\[
 A=x^2,\qquad C=[x,y],\qquad B=y,
\]
and filter by the number of occurrences of \(B\).
The pentagon for a leading part becomes a relation in the enveloping
algebra of the level-two cyclotomic hyperplane arrangement
\cite[Section~1.4]{Enriquez}, modulo terms of lower count.
The comparison uses a common ordering procedure: corrections
raise the count in the braid algebra and lower it in the arrangement
algebra.  Keeping only the common count-preserving terms transfers the
relation without any independence assumption on the target algebra.

Integral deletion operators on words then detect the leading part
after setting \(A=0\).  Their analytic antecedents are the
total-differential formulas for iterated integrals and hyperlogarithms
of Goncharov \cite[Theorem~2.1]{Goncharov} and
Panzer \cite[Lemma~3.3.30]{Panzer}.
The operators used here specialize the integral
algebraic differential operators of Hirose and Sato
\cite[Section~2]{HiroseSato}; their use with cyclotomic braid algebras
and ratios of sums and differences follows Hirose's construction
\cite[Sections~4.3 and~6]{Hirose}.  Here the arrangement relations are proved over the
integers before reduction, and the characteristic-two comparison is
proved directly.  A shift symmetry
\(B\mapsto B+s(C)\), followed by a Vandermonde argument, confines the
image to the ordinary free Lie algebra on
\[
 \ad_C^kB,\qquad k\ge1.
\]
These letters have exactly the weights predicted by the conjecture.
The ordering comparison and the integral deletion test provide
separate tools for extracting information from filtered braid
relations.

To attain the bound within the proof, we construct rational solutions with a nonzero
coefficient of \(x^{n-1}y\) in every odd weight \(n\ge3\).
Following Drinfeld's associator proof
\cite[Proposition~6.3]{Drinfeld}, we compare regularized braid holonomy
with the rule obtained by conjugating paths and negating the chord
generators.  We implement
this comparison in the parenthesized-chord formalism of Bar-Natan
\cite[Sections~2--3]{BarNatan}, with the required operations and
limits proved explicitly.  The logarithm of the comparison,
evaluated on a reassociation arrow, has the required nonzero
coefficient.  This uses the nonvanishing of a convergent integral,
and requires no arithmetic independence assertion about its values.

Finally, we work with saturated lattices over \(\Z_{(2)}\).
The characteristic-two bound applies to reductions of rational
solutions.  Independent Hall words in previously chosen generators,
together with a new depth-one value in odd weight, attain the bound
at each step.
Equality forces the next integral generator to occupy the required
leading filtration piece.  This proves generation and freeness
simultaneously.

\subsection{Organization and conventions}
Section~\ref{sec:algebra} establishes the integral algebra and a
special-derivation identity.  Sections~\ref{sec:degeneration}--\ref{sec:image}
prove the dimension bound.  Sections~\ref{sec:category}
and~\ref{sec:holonomy} construct the odd-weight values.
Section~\ref{sec:induction} completes the induction.

All Lie algebras in characteristic two are \emph{ordinary} Lie
algebras, with \([u,u]=0\).  Their enveloping algebras are ordinary
associative enveloping algebras.  We never identify all primitive
elements in characteristic two with ordinary Lie elements.
All completions are by weight, and all infinite constructions are
interpreted first in each finite weight quotient.
